\documentclass[11pt]{article}
\usepackage[margin=1in]{geometry}
\usepackage[T1]{fontenc}
\usepackage{lmodern}
\usepackage{microtype}
\usepackage{booktabs}
\usepackage{tabularx}
\usepackage{amsmath,amssymb}
\usepackage{enumitem}
\usepackage[hidelinks]{hyperref}
\usepackage{url}
\setlist{itemsep=2pt,topsep=4pt}
\newcommand{\code}[1]{\texttt{#1}}
\newcolumntype{R}{>{\raggedleft\arraybackslash}X}

\title{YonedaRepo: A Global-First Git Forge\\Built as a Code-as-Transactions Database}
\author{Matthew Long\\YonedaAI}
\date{Draft, 4 November 2026}

\begin{document}
\maketitle

\begin{abstract}
Coding agents change the load on a source forge: many writers, one trunk, and changes
that arrive seconds apart. We describe YonedaRepo, a Git-compatible forge that treats a
repository as a small transactional database while keeping Git's object model unchanged.
Every repository has one authority, a SQLite-backed Cloudflare Durable Object, that keeps
an append-only, hash-chained transaction log of ref transitions in the style of Delta
Lake. Commits are built by stateless compute and recorded in sharded objects; a landing
is a compare-and-swap of a branch plus one log entry, validated optimistically against
the entries written after the change's base. Changes may declare what they read, down to
a single symbol, and a landing whose reads were overwritten is refused. Path-disjoint
landings are batched into changesets and projected to ordinary Git commits afterwards.
The forge is global-first: requests run in the caller's region, new repositories are
placed near their creator, and verified immutable data is copied to readers' regions
behind epoch fences.
On the hosted development environment, batched landing takes 1,000 concurrent writers
on one branch to 1,000 of 1,000 landed and read back, at 404--491 landings/s in three
runs, where a polling design landed 4 of 1,000. Moving the request path into the
caller's region cut median landing latency from 205 to 53\,ms in Europe and from 382
to 33\,ms in Asia. We state the guarantees precisely: each landing is an atomic,
durable transaction in one repository, serializable over the reads a change declares.
Without declared reads, isolation is close to snapshot isolation. We report the gaps
that remain, including loaded read latency above our 1\,s target.
\end{abstract}

\section{Introduction}

A conventional forge assumes that people write code, that a pull request waits hours
for review, and that conflicts are rare enough to resolve by hand. Coding agents break
each assumption. They are fast, numerous, and often edit the same trunk at the same
time. When a thousand agents commit to one branch in the same few seconds, the forge's
job is no longer to store files; it is to order concurrent transactions and prove what
each one saw.

YonedaRepo is a forge built for that load. Its central decision is to treat source
control as a database problem without leaving Git. Git already gives an immutable,
content-addressed object store and atomic ref updates. What it does not give is a
transaction model across many concurrent writers: an order, a validation rule, a
durable record of what landed and why, and a defined isolation level. YonedaRepo adds
those around Git, never inside it.

\paragraph{Contributions.}
\begin{enumerate}
\item An architecture that keeps Git's object model as the foundation and places a
  per-repository transaction log, commit shards and stateless compute around it on
  serverless primitives (\S\ref{sec:arch}, \S\ref{sec:log}).
\item Optimistic concurrency at landing with declared read sets at file, directory or
  symbol granularity, plus batched path-disjoint changesets that remove the single
  branch's merge commit from the critical path (\S\ref{sec:landing}).
\item A global-first placement design: in-region execution, origin-placed authorities
  and epoch-fenced read copies, with the platform limits that bound it
  (\S\ref{sec:global}).
\item Current measurements of 1,000-writer bursts, sustained landing and regional
  latency on a hosted environment, with their run counts and failures
  (\S\ref{sec:eval}), and a precise statement of the transaction guarantees and their
  gaps (\S\ref{sec:guarantees}).
\end{enumerate}

The theory behind the design (the transaction model, the independence condition and a
Lean development of batch landing) appears in a companion series~\cite{long2026series}.
This paper is the systems report: what was built, how it behaves under load, and what
it does not yet guarantee.

\section{Design principles}
\label{sec:principles}

\begin{enumerate}
\item \textbf{Git's object model is the foundation.} Blobs, trees, commits and tags stay
  standard, immutable and content-addressed. There are no new object types, no extra
  headers and no rewritten history. Any Git client clones, fetches and pushes an
  ordinary repository.
\item \textbf{Transactions act on refs, never on objects.} The log records ref
  transitions and the paths each touched. Replaying it rebuilds the refs.
\item \textbf{One authority per repository.} A repository is the unit of serialization.
  Its authority's SQLite database holds the log, refs, approvals and grants; every
  other store is a rebuildable projection.
\item \textbf{No network I/O inside a transaction.} State, ordered events and the outbox
  commit together; anything that talks to the network runs before or after, driven by
  the outbox.
\item \textbf{Agents are clients.} Agents use stock Git, a command-line client and MCP.
  The platform does not orchestrate them. What an agent claims (that it read a file,
  that tests passed) is an assertion, checked against the exact commit that lands.
\item \textbf{Locality follows the request.} Work runs near the agent or person who
  asked, and data does not cross continents to answer them.
\end{enumerate}

\section{Architecture}
\label{sec:arch}

YonedaRepo runs entirely on Cloudflare. Rust owns the domain rules and durable state
transitions and is compiled to WebAssembly for Workers and Durable Objects; a thin
TypeScript layer adapts the platform SDKs. Table~\ref{tab:components} lists the
components.

\begin{table}[t]
\centering
\small
\begin{tabularx}{\linewidth}{@{}lXl@{}}
\toprule
Component & Role & Authoritative \\
\midrule
SourceRepository object & Per-repository authority: log, refs, approvals, grants (SQLite) & Yes \\
CommitShard objects (64) & Record commits built by stateless compute, chosen by commit id & For their records \\
LandCombiner object & Branch combiner and head holder; the authority's CAS decides & No \\
WorkspaceAuthority object & One per account: sessions, tokens, device logins & For accounts \\
Read-copy objects & Per-region copies of verified immutable trees and small blobs & No; epoch-fenced \\
R2 & Git objects (store of record), log segments, bundles, packs, LFS & Yes, for objects \\
D1 & Repository index and discovery & No; rebuildable \\
Queues, Workflows & Derived data, ref projection, imports (at least once) & No \\
\bottomrule
\end{tabularx}
\caption{Components. Authorities are Durable Objects with SQLite storage.}
\label{tab:components}
\end{table}

\paragraph{Commits without the authority.} Building a commit touches no ref. Stateless
compute builds the Git objects and records the commit in one of 64 commit shards,
chosen by the first byte of the commit id; a 1,000-writer burst places about 16 commits
on each. The authority adopts a commit only when it first needs it, at landing.

\paragraph{Objects and staging.} R2 is the store of record for Git objects, keyed
\path{platform/v1/<repo>/git/<id>} and never deleted. Because an R2 write takes on the
order of a second, a small commit is accepted once its authority durably holds every
object it built. The same transaction writes an object-outbox row; a projector writes
the objects to R2, confirms each write, and only then drops the staged copies. Reads fall
back to staged copies by id, and nothing is dropped unwritten. Large pushes are written as bundles (frames of kind, 20-byte id,
length and content) because a Worker request may have at most six R2 operations in
flight; one object per R2 write made a 12,000-file push take minutes.

\paragraph{Workers.} A front Worker serves the browser application and sign-in and
proxies the API to a platform Worker. Six regional API Workers (Section~\ref{sec:global})
share the same Durable Objects and storage and own no duplicates of them.

\section{The transaction log}
\label{sec:log}

\paragraph{Entries.} Each entry records one ref transition: its version, the base
version the change started from, a kind (commit, merge, revert, import, ref move, ref
delete or changeset), the ref name, old and new commit ids, the paths touched relative
to the first parent (with operation, old and new blob, and mode), the declared read set,
the actor, and the commit time. Validation requires the base to precede the version, a
fully qualified ref, no duplicate path, and a correct chain digest.

\paragraph{Hash chain and checkpoints.} An entry's chain digest is
\[ c_v = \mathrm{sha256}\bigl(c_{v-1} \,\|\, \text{entry}_v \text{ without its chain}\bigr), \]
streamed path by path so a large fast-forward does not have to be held as one value. Every 64
versions a checkpoint records the full ref map and its digest; the next entry carries
the digest of the checkpoint before it, as Delta Lake's version checksum
does~\cite{armbrust2020delta}. Replay checks contiguity, the chain, and each ref's old
value.

\paragraph{Reads at a version.} Reads take a version, a time or a commit, and answer
from the log's snapshot at that point.

\paragraph{Archive and compaction.} The log is sealed into segments of 1,024 versions,
written to R2 with the fingerprint of their content. Compaction writes and reads back a
segment, rebuilds it from the rows, requires identical fingerprints, then deletes the
rows and advances the anchor in one transaction; a rerun is idempotent. A read below the
floor answers with the names of the archived objects; the adapter fetches them and the
authority accepts them only with their recorded fingerprints, keeping network I/O out
of the authority.

\section{Concurrency control at landing}
\label{sec:landing}

\subsection{Optimistic validation}

A landing moves a branch by compare-and-swap and appends one entry. There is no merge
lock. The authority performs backward validation~\cite{kung1981occ} against the entries
on the target ref after the change's base, which are serialized before it:

\begin{description}
\item[Fast-forward] nothing moved the ref; swap directly.
\item[Rebase] later entries wrote only other paths; build a merge commit.
\item[Merge] later entries wrote some of the same files and nothing this change read;
  attempt a three-way content merge.
\item[Conflict] a later entry wrote something this change read, or moved the ref
  without path detail (a reset or delete cannot be proven independent).
\end{description}

Two paths overlap when they are equal or one is a directory prefix of the other. The
platform computes merge commits with its own Rust merge; the authority checks that a
merge commit's parents are the current head and the landed work, that its principal is
the platform, and that it changes no path the merged work did not.

\subsection{Declared reads and symbols}

A change may declare its read set with \code{git push -o yoneda.reads=\ldots}, through
the \code{yoneda push} client, or through agent hooks that record the files an agent
reads. An entry is a file, a directory prefix, or a symbol such as
\code{src/lib.rs\#Policy::backoff}. Symbol writes are derived per commit from both sides
of the diff; an unparsed file, or a commit whose symbols are not derived, counts as a
whole-file write, so an unknown never reads as disjoint. A later write then conflicts
with a symbol read only if it touched that symbol.

\subsection{Batches, changesets and the queue}

Landing one change per authority call does not scale to a thousand writers on one
branch. In batch mode, the deployed default, a combiner submits every waiting landing on
a branch in one call of up to 256. The authority re-checks each member's credentials
(scope, path grant, pause, revocation) and partitions the batch:
\begin{itemize}
\item path-disjoint members become \emph{changesets}: one indexed query each, no Git
  objects and no merge commit at landing time. A deterministic projector later builds
  one ordinary Git commit for them;
\item members that write the same files take the merge path;
\item a member that would read or write what an earlier member of the same batch wrote
  waits, which prevents write skew within a batch (Adya G2-item~\cite{adya2000isolation}).
\end{itemize}
A merge queue lets whoever holds a branch's advisory lease (10\,s) merge pending
independent landings into one octopus merge; a batch lands only if its base is still
the head, so a late builder loses the swap and a vanished one never blocks the branch.

\subsection{Checks and clearance}

A pure function decides whether a landing is cleared, held (a pending check or approval,
an unresolved comment, a full repository) or sent around (a failed required check, a
rejection, a stale read, a conflicted merge). Every caller (the page, the background
job, the land call and the batch) asks the same function, and the land call decides
again inside its own transaction with the expected head. After main moves, platform
checks run on the combined head; a failure proposes a revert commit.

\section{Global-first placement}
\label{sec:global}

\paragraph{Rule.} Data does not cross continents to answer a request. A request from
Europe is served by Europe's Workers, objects and storage, and likewise elsewhere.
Only the platform's trusted view of the caller's continent selects a region; headers,
cookies and request bodies cannot.

\paragraph{Execution and placement.} Six regional API Workers run in North America,
Europe, Asia-Pacific, Oceania, South America and Africa. A new repository's authority
and shards are created with a location hint taken from the creator's continent before
first access. Durable Objects never move, so existing authorities stay where they were
created and identifiers never change. A request that reaches the wrong cell is
forwarded: it can be answered slowly, never wrongly. There is no automatic write
failover; an unavailable region fails rather than start a second history.

\paragraph{Read copies.} Verified immutable data (trees and small blobs) is copied into
the reader's region: at most 16 objects and 4\,MiB per request, 1\,MiB per object, and
512\,MiB per store. Commits, refs, large files and access grants are never answered from
a copy. Deletion is fenced: a registry advances a publication epoch before a purge fans
out, each store keeps the highest epoch and deletes bytes in the same SQL transaction, a
late fill with an older epoch is rejected, and the purge job cannot acknowledge until
every enrolled region is fenced.

\paragraph{Platform limits.} The design is bounded by what the platform offers as of
October 2026~\cite{cfdurableobjects}. Location hints are best effort and some regions do
not host objects; jurisdictions exist for the EU, the US and FedRAMP but not for Asia;
queues have no location hint; D1 writes always go to the primary; and a container is not
guaranteed to run where its object does.

\section{Git compatibility}
\label{sec:git}

YonedaRepo serves smart HTTP and protocol v2 for fetch (including partial clone with
\code{blob:none} and shallow clone), receive-pack for push, Git LFS through presigned
storage URLs, and SSH through a tunnel to a gateway container that bridges to smart HTTP.
A pushed pack is verified before any ref moves: every introduced tree and file must be
in the pack or already in the repository, and each commit's bytes must hash to its id.
Objects are served only if reachable from a current branch or tag. Working branches move
by compare-and-swap on the value Git saw; \code{main} is protected, so a force push lands
like any other push (merge or conflict) and never rewrites it. A push holds at most
64\,MiB of pack in memory and 40,000 objects; larger files stream to storage.

\section{Evaluation}
\label{sec:eval}

\paragraph{Method.} All measurements below are from the hosted development environment
unless marked local. The load-testing rules require at least three runs per size and
server-side timing; where a result has fewer runs we say so. In every burst each writer
commits one new file on \code{main} and lands it, with no conflicting writes and no
declared reads, from a benchmark Worker in \code{aws:us-east-1}. The bursts therefore
measure throughput and latency of the landing path, not its conflict handling, and are
not evidence for the isolation guarantees in Section~\ref{sec:guarantees}.

\subsection{1,000 writers on one branch}

\paragraph{Landing designs.} Table~\ref{tab:modes} compares three landing designs at
one revision with one run each. Polling the authority per landing collapsed (994 of
1,000 refused as overloaded); a per-landing combiner timed out 615 landings; changesets
landed all 1,000. The comparison varies the whole landing path, so it does not isolate
the cost of building Git objects.

\begin{table}[t]
\centering
\small
\begin{tabular}{@{}lrrr@{}}
\toprule
Landing design & Landed / 1,000 & Commit p95 (ms) & Land p95 (ms) \\
\midrule
Polling, Git built per landing & 4 & 2,221 & 104,034 \\
Combiner, Git built per landing & 385 & 2,704 & 64,756 \\
Changesets & 1,000 & 1,698 & 18,728 \\
\bottomrule
\end{tabular}
\caption{Landing designs, 1,000 writers, one run each (revision \code{dd6c9a8}). Client-side p95.}
\label{tab:modes}
\end{table}

\paragraph{Batch mode.} Table~\ref{tab:batch} shows batch mode, the deployed default.
Across six runs no commit, landing or read was refused and all 1,000 reads returned the
correct content. Separate tokens per writer landed 404--491 changes per second over the
burst; one shared login for all writers landed 235--284.
Main's Git commit trailed the log by 3.2--5.4\,s at p95, because changesets are
projected after landing. Three later runs at revision \code{19d58ad} (27 October)
landed 1,000 of 1,000 each at 259--303 landings/s with land p95 of 2.7--3.0\,s.

\begin{table}[t]
\centering
\small
\begin{tabular}{@{}lrrrr@{}}
\toprule
Writers' credentials & Runs & Landings/s & Commit p95 (ms) & Land p95 (ms) \\
\midrule
One token per writer & 3 & 404--491 & 714--1,673 & 1,212--1,865 \\
One shared login & 3 & 235--284 & 1,140--1,614 & 2,167--3,027 \\
\bottomrule
\end{tabular}
\caption{Batch mode, 1,000 writers, revision \code{fef94e7}; 1,000 of 1,000 landed in
every run. Ranges are across runs; client-side p95.}
\label{tab:batch}
\end{table}

\paragraph{Current revision.} Table~\ref{tab:current} gives server-side timing for three
runs on 3 November 2026 (revision \code{c8d6d9ae}), with 1,000 temporary-token writers
each. All 3,000 commits, landings and exact-content readbacks passed. Client read p95
across runs was 1,569 / 1,703 / 2,807\,ms (min / median / max; not a pooled p95). Three
runs the same day at revision \code{b38706d6} also passed 3,000 of 3,000, with read p95
of 1,732 / 2,644 / 2,796\,ms. \textbf{Correctness held; subsecond reads under load did
not}: our budget is a read p95 under 1\,s during any burst. Fetching the head
accounts for most of the server's read time (p95 969--2,309\,ms of 1,436--2,735\,ms);
reads served from the edge took 35--48\,ms at p95.

\begin{table}[t]
\centering
\small
\begin{tabular}{@{}lrrrr@{}}
\toprule
Run & Commit (server) & Land (server) & Read (client) & Head fetch (server) \\
\midrule
1 & 2,981 / 3,218 / 3,356 & 679 / 1,349 / 1,412 & 1,524 / 1,703 / 7,578 & 819 / 1,528 / 1,623 \\
2 & 1,862 / 2,228 / 3,280 & 501 / 755 / 2,638 & 2,553 / 2,807 / 4,948 & 1,650 / 2,309 / 2,614 \\
3 & 1,842 / 2,179 / 2,681 & 833 / 1,148 / 1,495 & 1,382 / 1,569 / 4,386 & 760 / 969 / 1,469 \\
\bottomrule
\end{tabular}
\caption{1,000-writer bursts at revision \code{c8d6d9ae}, 3 November 2026. Each cell is
p50 / p95 / max in ms over 1,000 samples.}
\label{tab:current}
\end{table}

\paragraph{Between revisions.} Nineteen further 1,000-writer runs from 1 to 4 November,
one per revision, all landed and read back 1,000 of 1,000 with no overload refusals, at
137.9--557.7 landings/s and server land p95 of 593--2,986\,ms. With one run per revision
they form a time series, not per-revision results, and we do not attribute the spread
to any change. Their evidence files are not yet published with the repository.

\paragraph{What the bursts do not show.} Three bursts of 1,000 are not 3,000
simultaneous writers. They use API writers, not real agents or native Git pushes, and
none ran in production. More than 1,000 writers has not been measured.

\subsection{Sustained landing and write skew}

An open-loop test on one repository (27 October, one run per step of two minutes) held
100 landings/s offered (99.6/s achieved, p95 2.6\,s, 0.02\% refused) and missed 200/s
(123/s achieved, 17\% refused), placing the knee between the two. At a later revision
(\code{75e7752}), 200/s offered for 60\,s achieved 193.7/s with none refused and p95 3.4\,s
(one run).

A write-skew workload of 200 writers, three runs, found 0 serializability violations
among 6,441--8,001 log entries checked; 114--127 changes landed with declared reads.

\subsection{Regional latency}

Table~\ref{tab:region} compares median client latency with the request path in North
America (today) against the same path placed in the client's region, three runs each on
1 November 2026 with zero failures. The path ran in Amsterdam for Europe and Singapore
for Asia in all runs. Reads, lists, commits and landings, which the repository's own
objects answer, fell to 4--53\,ms in region. Search and repository creation became slower, because D1 and R2 stayed in
North America and moved further from the code; a personal-token read stays bounded by
D1. For reference, North America today answers a read in 41\,ms and a landing in
113\,ms.

\begin{table}[t]
\centering
\small
\begin{tabular}{@{}lrr@{}}
\toprule
Operation & Europe & Asia \\
\midrule
Read (repository token) & 162 $\to$ 22 & 275 $\to$ 15 \\
List & 144 $\to$ 11 & 256 $\to$ 4 \\
Commit & 224 $\to$ 51 & 342 $\to$ 41 \\
Land & 205 $\to$ 53 & 382 $\to$ 33 \\
Sign-in & 244 $\to$ 123 & 398 $\to$ 76 \\
Push & 444 $\to$ 251 & 763 $\to$ 426 \\
Search & 496 $\to$ 566 & 670 $\to$ 1,142 \\
Repository create & 2,439 $\to$ 2,959 & 2,361 $\to$ 3,814 \\
\bottomrule
\end{tabular}
\caption{Median client latency (ms), request path in North America $\to$ in region.
Clients in \code{aws:eu-central-1} and \code{aws:ap-southeast-1}; 3 runs.}
\label{tab:region}
\end{table}

On 3 November the six regional API Workers were each observed executing in their
region (IAD, FRA, SIN, SYD, GRU, CPT). Every test caller reported a US origin, so this
establishes execution placement, not the experience of a European or Asian user.

\subsection{Storage}

A one-file landing costs about 4\,KB in the authority once projected (1,024 landings in
four batches), so a repository reaches its 9.4\,GB internal stop after about 2.3 million
such landings. A spike on one object found writes refused at about 10.7\,GB; deleting
950\,MB in one transaction reset the object, which then missed two 120-second requests,
while a 19\,MB delete was fine. Planned history shards seal old history into read-only
objects before that limit; they are not built.

\subsection{Correctness testing}

An independent reference model of a repository, with no dependencies of any kind, is
the oracle for seeded random operation sequences (16 cases of 80 steps by default) that
must reach 22 named outcomes. After each step every outcome, merge plan, ref, version
and file is compared; at the end real Git checks the objects with
\code{fsck --strict} and the merge bases. Contract tests pin dirty write, atomic
multi-ref batches, lost update in both directions and stale reads, and one test records
the known write-skew gap described next.

\section{Guarantees and gaps}
\label{sec:guarantees}

\paragraph{Boundary.} A transaction is one landing: one ref transition, or one atomic
batch of working-branch moves, in one repository. Its commit point is the SQLite
transaction that moves the ref and appends the log entry. There are no transactions
across repositories.

\paragraph{What holds.}
\begin{itemize}
\item \emph{Atomicity}: the ref swap and the log entry commit together or not at all; an
  atomic push with one stale branch moves nothing.
\item \emph{Durability}: the reply waits until the authority's transaction is durable,
  and staged objects remain in the authority until storage confirms each write.
\item \emph{Isolation}: Table~\ref{tab:isolation}. With declared reads, landing is
  serializable over the declared read set. Without them, a landing is checked only for
  writes to the same files, which are merged three ways: lost updates are prevented but
  write skew~\cite{berenson1995critique} through undeclared reads can land.
\item \emph{Consistency}: required checks and approvals are judged on the exact head and
  decided again inside the land transaction.
\end{itemize}

\begin{table}[t]
\centering
\small
\begin{tabularx}{\linewidth}{@{}lX@{}}
\toprule
Path & Isolation \\
\midrule
Explicit landing, no declared reads & Close to snapshot isolation: no lost update; write skew possible \\
Explicit landing, declared reads & Serializable over the declared read set \\
Batch (combiner) & Serializable over declared reads; a member with none counts its writes as reads \\
Land when cleared & As batch, plus separation in queue order \\
\bottomrule
\end{tabularx}
\caption{Isolation by landing path, as built.}
\label{tab:isolation}
\end{table}

\paragraph{Gaps, each with a planned fix.}
\begin{enumerate}
\item \emph{Checks after merges.} A platform merge commit lands under the gates of the
  work it merges; checks on the merged tree run after main moves. Planned: hold the merge
  until its checks pass, then land it with the expected head.
\item \emph{Merge content.} Verification checks a merge commit's shape (parents,
  principal, paths), not that its tree is the correct three-way merge. Planned: recompute
  the merge in the authority and refuse a differing tree.
\item \emph{Undeclared reads.} Planned: a serializable policy in which a landing with no
  declared reads counts its writes as reads, as batches already do.
\item \emph{Read coverage.} Declared reads are agent claims. In one experiment the hooks
  missed files read through shell commands; serializability is only as strong as the
  read set.
\item \emph{Failure model.} Isolate crash, object restart, storage failure and lost
  acknowledgements are not yet each covered by a named test.
\end{enumerate}

\section{Related work}
\label{sec:related}

\paragraph{Transactional Git storage.} Gitaly's transaction manager gives ACID writes
to refs, objects and hooks with serializable snapshot isolation and a write-ahead
log~\cite{gitaly2024tx}. Its conflicts are defined on repository state; YonedaRepo's
landing validation is defined on what a change read in the source tree. Piper serves
Google's monorepo from Bigtable and later Spanner~\cite{potvin2016piper}; it is not
stock-Git compatible. Sapling's server performs server-side rebase
(pushrebase)~\cite{sapling}, which corresponds to our rebase outcome for disjoint
writes.

\paragraph{Keeping trunk green.} SubmitQueue orders changes speculatively and treats
changes as conflicting when they touch overlapping build packages
\cite{ananthanarayanan2019submitqueue,juloori2025ci}; merge trains do the same with
pipelines~\cite{gitlabmergetrains}. These validate by rebuilding. YonedaRepo validates
by read set at landing and checks the combined result afterwards.

\paragraph{Transaction logs over object stores.} Delta Lake keeps an ordered log with
optimistic commits and checkpoints over cloud object storage~\cite{armbrust2020delta};
our log follows its shape for refs instead of table files. Dolt~\cite{dolt} and
Git4Data~\cite{gou2026git4data} bring Git-style branching to relational data; Jujutsu
keeps an operation log on the client~\cite{jujutsu}.

\paragraph{Concurrency for agents.} CoAgent fixes a serialization order for agents
sharing state such as a Git tree and repairs conflicting writes in middleware
\cite{lyu2026coagent}. Cognitive Serializability asks that an agent's commit revalidate
its data, evidence, policy and authority~\cite{he2026cognitive}, close in spirit to our
declared reads and exact-head clearance, though not applied to source control.
MetaGPT~\cite{hong2023metagpt} and CAID~\cite{caid2026} coordinate agents through
workflows and isolated worktrees; YonedaRepo deliberately leaves orchestration to the
clients.

To our knowledge, no prior forge validates landings against declared read sets at
symbol granularity while serving stock Git. Our search was not exhaustive.

\section{Conclusion}

Treating a repository as a database, with Git as its immutable object layer and a
per-repository log as its transaction record, lets a forge take a thousand concurrent
writers on one branch and land all of them, and lets it run near its users without a
global database. The guarantees are real but scoped: atomic and durable per landing,
serializable over declared reads. The remaining work is to close the isolation and
consistency gaps listed above, measure beyond a thousand writers, and bring loaded read
latency under one second.

\bibliographystyle{plain}
\bibliography{refs}

\end{document}
